% ============================================================
%  Chapter 3 — How Robots Sense the World
%  Robot World: Meet the Machines That Help Us
% ============================================================

\chapter{How Robots Sense the World}

\chapteropening{%
  You have eyes to see, ears to hear, and skin to feel.
  Robots have sensors \textemdash{} electronic eyes, ears, and fingertips
  that help them understand what is happening around them.
}
\chapterbadge{Sensors: A Robot's Window to the World}

% --- Introduction ---
\section*{Why Sensing Matters}

Imagine trying to walk through your house with your eyes closed and your
ears plugged. You would bump into everything! Robots without sensors would
do the same.

\character{Dot}{Sensors are the first step in the sense\,--\,think\,--\,act
  loop. Without them, a robot is just a lump of metal.}

% --- Cross-book callback (added per review: strengthen series continuity) ---
\begin{tcolorbox}[enhanced,colback=SunYellow!12,colframe=SunYellow!70!black,arc=6pt,boxrule=1pt,
  fonttitle=\bfseries, title={\faRetweet\quad Remember from Book 2?}]
  In \emph{Inside a Smartphone}, you met Sensa \textemdash{} the hidden sensors inside a phone:
  the \textbf{accelerometer} (movement), \textbf{gyroscope} (tilt), \textbf{GPS} (location),
  and \textbf{camera} (sight). Robot sensors work on the same principles, just bigger, tougher,
  and sometimes in combinations a phone would never need!
\end{tcolorbox}

\sectionrule

% --- Types of sensors ---
\section*{Types of Robot Sensors}

\subsection*{Cameras \textemdash{} Robot Eyes}

\term{Cameras} capture pictures many times per second. A robot's computer analyses
these images to recognise objects, read signs, or follow a line on the floor.

\begin{funfact}
  Self-driving cars use up to eight \term{cameras} at once, giving them a
  360-degree view \textemdash{} humans can only see about 180 degrees!
\end{funfact}


\begin{center}
\begin{tabularx}{0.88\textwidth}{|X|X|X|}
  \hline
  \textbf{Everyday Device} & \textbf{Sensor Used} & \textbf{What It Detects} \\
  \hline
  Automatic door & Infrared / motion & A person approaching \\
  \hline
  Smartphone & Accelerometer & Tilt (screen rotation) \\
  \hline
  Smoke alarm & Optical / ionisation & Smoke particles in air \\
  \hline
  Reversing car & Ultrasonic & Objects behind the car \\
  \hline
  Fitness tracker & Heart-rate sensor & Your pulse \\
  \hline
\end{tabularx}
\end{center}

\sectionrule

\subsection*{Lidar \textemdash{} Laser Eyes}

\term{Lidar} fires invisible \term{laser} beams in all directions and measures
how long each beam takes to bounce back. This creates a \term{3D map} of the
surroundings, even in the dark.

\subsection*{Ultrasonic Sensors \textemdash{} Echo Location}

Just like a bat, these \term{ultrasonic} sensors send out a pulse of sound and listen for the
echo. The time the echo takes to return tells the robot how far away an
object is.

\begin{picturethis}
  Stand in a large empty room and clap your hands. You hear the echo
  bounce off the walls. An ultrasonic sensor does the same thing, but
  millions of times faster and with sounds too high for your ears to hear.
\end{picturethis}

\subsection*{Touch and Pressure Sensors}

Some robots have sensitive skin. Pressure sensors tell them how hard
they are gripping something \textemdash{} essential for a robot handing you
a glass of water without crushing it!

\subsection*{Temperature Sensors}

Useful for robots that cook, fight fires, or explore volcanoes.
They measure heat without needing to touch the surface.

\subsection*{Gyroscopes and Accelerometers}

These sensors tell a robot which way is up, how fast it is moving,
and whether it is tilting. Drones would tumble out of the sky without them.

\character{Whirr}{My gyroscope keeps me level, even when the wind blows.
  Without it I would flip upside down in seconds!}

\sectionrule

% --- How sensors work together ---
\section*{Sensor Fusion \textemdash{} Combining Superpowers}

One sensor alone can be fooled. A camera cannot see in the dark. Lidar
cannot read colours. So engineers combine several sensors and let the
computer brain mix all the data together. This is called \term{sensor fusion}.

\begin{bigidea}
  Sensor fusion means combining information from many sensors
  to build a more complete picture of the world \textemdash{}
  just as you use sight \emph{and} hearing \emph{and} touch at the same time.
\end{bigidea}

\sectionrule

% --- Experiment ---
\begin{experiment}
  \textbf{Test your own sensor fusion.}
  \begin{enumerate}
    \item Close your eyes and try to touch your nose with one finger.
    \item Now spin around three times, \emph{then} try again.
    \item Notice how much harder it is when your balance sensor (inner ear)
      is confused \textemdash{} even though your touch sensor still works.
  \end{enumerate}
  \textit{Robots face the same challenge: if one sensor gives bad data,
  the brain must rely on the others.}
\end{experiment}

\sectionrule

% --- Diagram ---
\begin{diagram}[Common robot sensors and what they detect.]
  \begin{tikzpicture}[>=Stealth, node distance=1.6cm]
    \node[draw=RoboGreen!80!black,fill=RoboGreen!15,rounded corners=8pt,
          minimum width=2.2cm,minimum height=2.8cm,line width=1.2pt,
          font=\bfseries\small,align=center] (robot) at (0,0) {Robot\\Brain};
    % Sensors
    \node[draw=WarmOrange,fill=WarmOrange!10,rounded corners,
          minimum width=2.8cm,font=\small,right=2.5cm of robot,yshift=2.4cm] (cam) {Camera: light \& colour};
    \node[draw=SteelBlue,fill=SteelBlue!10,rounded corners,
          minimum width=2.8cm,font=\small,right=2.5cm of robot,yshift=1.0cm] (lid) {Lidar: distance (laser)};
    \node[draw=SoftPurple,fill=SoftPurple!10,rounded corners,
          minimum width=2.8cm,font=\small,right=2.5cm of robot,yshift=-0.4cm] (ult) {Ultrasonic: distance (sound)};
    \node[draw=CoralRed,fill=CoralRed!10,rounded corners,
          minimum width=2.8cm,font=\small,right=2.5cm of robot,yshift=-1.8cm] (tch) {Touch: pressure};
    \node[draw=SunYellow!80!black,fill=SunYellow!10,rounded corners,
          minimum width=2.8cm,font=\small,left=2.5cm of robot,yshift=0.8cm] (gyr) {Gyroscope: tilt \& spin};
    \node[draw=LeafGreen,fill=LeafGreen!10,rounded corners,
          minimum width=2.8cm,font=\small,left=2.5cm of robot,yshift=-0.8cm] (tmp) {Thermometer: heat};
    % Arrows
    \draw[->,RoboGreen!60!black] (cam.west) -- (robot.east|-cam.west);
    \draw[->,RoboGreen!60!black] (lid.west) -- (robot.east|-lid.west);
    \draw[->,RoboGreen!60!black] (ult.west) -- (robot.east|-ult.west);
    \draw[->,RoboGreen!60!black] (tch.west) -- (robot.east|-tch.west);
    \draw[->,RoboGreen!60!black] (gyr.east) -- (robot.west|-gyr.east);
    \draw[->,RoboGreen!60!black] (tmp.east) -- (robot.west|-tmp.east);
  \end{tikzpicture}
\end{diagram}

\sectionrule


\begin{quickcheck}
  \begin{enumerate}
    \item What does a lidar sensor measure, and how?
    \item Why is sensor fusion better than using one sensor alone?
    \item Which sensor would a robot need to detect a hot surface?
  \end{enumerate}
\end{quickcheck}

\sectionrule
% --- New Words ---
\begin{newwords}
  \begin{description}[labelwidth=3.8cm, leftmargin=4.0cm]
    \item[Sensor] A device that measures something in the environment.
    \item[Lidar] Light Detection and Ranging \textemdash{} uses lasers to map distances.
    \item[Ultrasonic] Using sound waves above human hearing to detect objects.
    \item[Sensor fusion] Combining data from multiple sensors for a better picture.
    \item[Gyroscope] A sensor that measures rotation and tilt.
    \item[Accelerometer] A sensor that measures changes in speed or direction.
  \end{description}
\end{newwords}
